\documentclass[11pt]{article}
\input{packages.tex}
\DeclareMathOperator{\sign}{sign}
\DeclareMathOperator{\Imm}{Im}

% Настройка колонтитулов
\pagestyle{fancy}
\fancyhf{} % Очистить все поля
\fancyhead[L]{Мы весёлые медузы, мы похожи на арбузы}
\fancyhead[R]{\href{https://t.me/sol_m_07}{Соля}}
\fancyfoot[C]{\thepage}               % Справа внизу

\setlength{\headheight}{13.6pt}

\hypersetup{
    colorlinks=true,
    linkcolor=blue,
    citecolor=blue,
    filecolor=blue,
    urlcolor=blue,
}

\begin{document}

\begin{titlepage}
    \centering
    \vspace*{\fill}  % заполняет пространство сверху

    {\Huge Лекция 1. ЛЗ и ЛНЗ строк/столбцов матрицы. Ранг матрицы (строчный, столбцовый, минорный)
    \par}

    \vspace*{\fill}  % заполняет пространство снизу
\end{titlepage}

\setcounter{page}{1}

\section{ЛЗ и ЛНЗ строк/столбцов матрицы}

Вспоминая начало предыдущего семестра, по аналогии с наборами векторов начнём 
рассматривать наборы матриц. Тогда линейной комбинацией матриц назовём: 
\[\alpha_1A_1 + \alpha_2A_2 + ... + \alpha_kA_k\]
Точно так же будем называть набор линейно независимым (ЛНЗ), если только тривиальная 
комбинация обращается в ноль. Иначе - линейно зависимым (ЛЗ). 

\begin{remark}
    Все прошлые теоремы (в том числе критерии ЛЗ и ЛНЗ) выполняются и для матриц.
\end{remark}

\begin{theorem}
    Для любой квадратной матрицы выполнено 
    \[detA \neq 0 \Longleftrightarrow \text{строки ЛНЗ} \Longleftrightarrow \text{столбцы ЛНЗ}\]
\end{theorem}

\begin{theorem}
    Для любой квадратной матрицы следующие условия эквивалентны:
    \begin{enumerate}
        \item $detA = 0$
        \item $\exists c \neq O: cA = O$, с - строка
        \item $\exists c \neq O: Ac = O$, с - столбец
    \end{enumerate}
\end{theorem}

\begin{theorem}
    Пусть $A$ и $B$ - квадратные матрицы одного размера. Тогда если $A$ - вырожденная, 
    то и $AB$ - вырожденная.
\end{theorem}

\begin{proof}
    Из предыдущей теоремы $detA = 0 \Longleftrightarrow \exists c \neq O: cA = O$. 
    Тогда из ассоциативности получаем \[c(AB) = (cA)B = OB = O\] 
    Иначе можно было воспользоваться теоремой о том, что $det(AB) = detA \cdot detB$
\end{proof}

\begin{theorem}
    При совершении элементарных операций ЛНЗ (ЛЗ) строки переходят в ЛНЗ (ЛЗ).
\end{theorem}

\begin{proof}
    \textbf{Шаг 1} (ЛНЗ $\rightarrow $ ЛНЗ). Пусть матрица $A$ состоит из ЛНЗ строк 
    $a_1, a_2, ..., a_n$. Необходимо доказать, что при умножении строки на число и 
    при сложении строк набор останется ЛНЗ 
    (основные элементарные операции). Первое очевидно верно, так как просто в линейной 
    комбинации получим перед каким-то $a_i$ не $\alpha_i$, а $k \cdot \alpha_i$. 
    Рассмотрим случай со сложением строк. Без ограничения общности посмотрим на 
    следующий переход: \[\begin{pmatrix} a_1 \\ a_2 \\ ... \\ a_n \end{pmatrix} 
    \longrightarrow \begin{pmatrix} a_1 + a_2 \\ a_2 \\ ... \\ a_n \end{pmatrix}\]
    Поймём, что \[\alpha_1(a_1 + a_2) + \alpha_2a_2 + ... + \alpha_na_n = 0 
    \Longleftrightarrow \alpha_1a_1 + (\alpha_1 + \alpha_2)a_2 + ... + \alpha_na_n = 0\] 
    Но ведь $a_1, a_2, ..., a_n$ - ЛНЗ набор. Значит только тривиальная комбинация 
    обращается в ноль. \[\begin{cases} \alpha_1 = 0 \\ \alpha_1 + \alpha_2 = 0 \\ 
    ... \\ \alpha_n = 0 \end{cases} \implies \alpha_2 = 0 \implies \text{после ЭО строки 
    остались ЛНЗ}\]
    \textbf{Шаг 2} (ЛЗ $\rightarrow $ ЛЗ). Пусть матрица $A$ состоит из ЛЗ строк 
    $a_1, a_2, ..., a_n$. Тогда от противного предположим, что после ЭО набор 
    становится ЛНЗ. Так как ЭО обратимы, то вернуться к исходному набору можно через 
    некоторое количество ЭО. Но по выше доказанному получаем, что набор теперь ЛНЗ - 
    противоречие, так как он был ЛЗ.
\end{proof}

\begin{remark}
    Аналогичное утверждение справедливо и для столбцов.
\end{remark}

\begin{theorem}
    При ЭО со столбцами линейная зависимость строк сохраняется с теми же 
    коэффициентами. Аналогично и в обратную сторону.
\end{theorem}

\begin{proof}
    Пусть $S$ - матрица элементарных операций со столбцами. Из линейной зависимости 
    $A$ получаем, что $\exists c \neq O: cA = O$. Тогда \[c(AS) = (cA)S = OS = O\] 
    Значит линейная зависимость сохраняется после ЭО с теми же коэффициентами.
\end{proof}

\section{Ранг матрицы}

\begin{definition}
    Пусть в $A \ \exists$ ЛНЗ набор из $r$ строк, а любой набор из $k > r$ строк ЛЗ. 
    Тогда $r$ - строчный ранг матрицы. Обозначается $r_1(A)$. $A = O \implies r_1(A) 
    := 0$. Эти $r$ строк назовём базисными.
\end{definition}

\begin{theorem}
    Любая строка матрицы - линейная комбинация базисных строк. 
\end{theorem}

\begin{proof}
    Если строка базисная, то её берём с коэффициентом 1, а остальные - с нулями. Пусть 
    строка $\widetilde{a}$ не базисная, а $a_1, a_2, ..., a_r$ - базисные. Из определения 
    ранга получаем, что $a_1, a_2, ..., a_r, \widetilde{a}$ - ЛЗ набор. Значит 
    существует нетривиальная комбинация, обращающаяся в ноль: \[\alpha_1a_1 + ... + 
    \alpha_ra_r + \alpha_{r + 1}\widetilde{a} = 0\] Приэтом $\alpha_{r + 1} \neq 0$, так 
    как иначе получим $a_1, a_2, ..., a_r$ - ЛЗ набор. Значит можем поделить на 
    $\alpha_{r + 1}$ и выразить $\widetilde{a}$ через базисные векторы (просто перенести 
    их вправо).
\end{proof}

\begin{definition}
    Пусть в $A \ \exists$ ЛНЗ набор из $r$ столбцов, а любой набор из $k > r$ столбцов 
    ЛЗ. Тогда $r$ - столбцовый ранг матрицы. Обозначается $r_2(A)$. $A = O \implies 
    r_2(A) := 0$. Эти $r$ столбцов назовём базисными.
\end{definition}

\begin{remark}
    Для базисных столбцов также выполняется теорема (подобно 2.1). Доказательство аналогично.
\end{remark}

\begin{theorem}
    ЭО со строками и столбцами не меняют ни строчного, ни столбцового рангов. 
\end{theorem}

\begin{proof}
    Доказательство следует из теоремы о сохранении ЛЗ при ЭО.
\end{proof}

\begin{theorem}
    Пусть $A_1 \subset A$. Тогда \[r_1(A_1) \leq r_1(A)\] \[r_2(A_1) \leq r_2(A)\]
\end{theorem}

\begin{proof}
    Если в $A$ есть ЛЗ строки, то они (их фрагменты) будут ЛЗ и в $A_1$, притом с теми 
    же коэффициентами.
\end{proof}

\begin{remark}
    Верен не только факт из доказательства, но и противоположный. Если в $A_1$ строки 
    ЛНЗ, то и в $A$ они ЛНЗ. Для доказательства достаточно предположить противное и 
    получить, что ЛНЗ перешёл в ЛЗ, а при обратном переходе остался ЛЗ по предыдущей 
    теореме. 
\end{remark}

\begin{theorem}
    (1) $\forall n + 1$ столбцов высоты $n$ ЛЗ \\
    (2) $\forall n + 1$ строк длины $n$ ЛЗ
\end{theorem}

\begin{proof}
    (1) $a_1, a_2, ..., a_n, a_{n + 1}$ - столбцы высоты $n$ \\
    $a_1, a_2, ..., a_n$ - ЛЗ $\implies$ весь набор ЛЗ \\
    $a_1, a_2, ..., a_n$ - ЛНЗ. Тогда составим невырожденную матрицу из этих столбцов. \\
    \[\forall b \ \exists x: Ax = b \implies a_{n + 1} \ \text{может быть представлен как 
    линейная комбинация остальных}\] По критерию линейной зависимости набор $a_1, a_2, 
    ..., a_n, a_{n + 1}$ ЛЗ. \\
    (2) Транспонируем матрицу и ссылаемся на первый пункт.
\end{proof}

\begin{theorem}
    $\forall A \hookrightarrow r_1(A) = r_2(A)$
\end{theorem}

\begin{proof}
    $A_{m \times n}$
    \begin{figure}[h]
    \centering
    \includegraphics[width=0.5\linewidth]{pictures/Начальная матрица.png}
    \end{figure} \\
    Занулим не базисные строки с помощью ЭО, так как каждая из них выражается через 
    базисные (также можно сгруппировать базисные строки и столбцы для удобства). 
    Новую матрицу назовём $A'$. Так как мы просто провели некоторые ЭО, то ранги не 
    изменились: \[r_1(A) = r_1(A') \text{ и } r_2(A) = r_2(A')\]
    Далее образуем матрицу $D$ на пересечении базисных строк и столбцов.
    \begin{figure}[h]
    \centering
    \includegraphics[width=0.6\linewidth]{pictures/Изменённая матрица.png}
    \end{figure} \\
    Так как все элементы базисных столбцов, кроме тех, которые входят в $D$, нулевые, то 
    ЛНЗ этих столбцов говорит о том, что эти же столбцы в $D$ тоже ЛНЗ. Получаем, что 
    в $D$ $r_2(A)$ ЛНЗ столбцов и $r_1(A)$ строк (нам неважно, какие они). Буквально 
    по предыдущей теореме тогда верно, что \\ $r_2(A) \leq r_1(A)$, так как иначе число 
    столбцов превысит высоту столбца и мы получим ЛЗ набор. После транспонирования 
    получим обратную ситуацию на ЛНЗ строки. Тогда верно неравенство $r_1(A) \leq r_2(A)$.
    \\ Итого: \[\begin{cases} r_2(A) \leq r_1(A) \\ r_1(A) \leq r_2(A) \end{cases} 
    \implies r_1(A) = r_2(A)\]
\end{proof}

\begin{lemma}[Следствие]
    На пересечении базисных строк и столбцов стоит невырожденная матрица.
\end{lemma}

\begin{definition}[Теорема-определение]
    Если $r_1(A) = r_2(A) = r$, то $\exists$ матрица размера $r$ с $det \neq 0$. 
    Приэтом любая подматрица с большим порядком, чем $r$, будет вырожденной. Порядок 
    этой матрицы называют минорным рангом.
\end{definition}

\begin{theorem}
    В любой матрице \[r_1(A) = r_2(A) = \text{минорный ранг}\]
\end{theorem}

\begin{definition}
    Все эти ранги, раз уж они равны, будем называть рангом матрицы и обозначать $rgA$.
\end{definition} \leavevmode
\\
\\
Для ценителей левой верхней фразы прикрепляю:
\\
\begin{figure}[h]
\centering

\includegraphics[width=0.9\linewidth]{pictures/Медузы.jpg}
\end{figure}

\end{document}